%% =============================================================================
%% manuscript.tex — The American Mathematical Monthly
%% Title:    The continuous-coefficient Jensen equation:
%%           A note on three vestigial regularity hypotheses
%% Author:   [AUTHOR]
%% Source:   satellites/o3-maa/latex/manuscript.tex
%% Date:     2026-06-08
%%
%% Compilation (WSL Debian, TeX Live 2024):
%%   pdflatex manuscript
%%   bibtex   manuscript
%%   pdflatex manuscript
%%   pdflatex manuscript
%%
%% Class: MAA Monthly template (article + maa-monthly.sty).
%%
%% DOUBLE-BLIND SUBMISSION (Monthly policy):
%%   This file compiles TWO documents in sequence:
%%   (1) Title page with author details (page 1).
%%   (2) Anonymous manuscript body (pages 2+).
%%   To produce the ANONYMOUS-ONLY PDF, comment out the title-page
%%   block between the markers below and recompile.
%%
%% maa-monthly.sty already loads the following -- do NOT re-add them:
%%   times, pifont, graphicx, color,
%%   amsmath, amsthm, amsfonts, amssymb, url.
%% =============================================================================

\documentclass{article}
\usepackage{maa-monthly}

%% --- Additional packages not in maa-monthly.sty ---
\usepackage[utf8]{inputenc}     % UTF-8 source encoding
\usepackage[T1]{fontenc}        % Type-1 output encoding
\usepackage{microtype}          % improved microtypography (safe with Times)
\usepackage{booktabs}           % \toprule, \midrule, \bottomrule
\usepackage{array}              % extended column specifiers in tables
\usepackage{tabularx}           % X column type for proportional-width tables

%% --- Theorem environments (following the Monthly template exactly) ---
\theoremstyle{plain}
\newtheorem{theorem}{Theorem}
\newtheorem{proposition}[theorem]{Proposition}
\newtheorem{lemma}[theorem]{Lemma}
\newtheorem{corollary}[theorem]{Corollary}

\theoremstyle{definition}
\newtheorem{definition}[theorem]{Definition}
\newtheorem{remark}[theorem]{Remark}

%% --- Math abbreviations ---
\newcommand{\R}{\mathbb{R}}
\newcommand{\Q}{\mathbb{Q}}
\newcommand{\E}{\mathbb{E}}

%% =============================================================================
\begin{document}

%% ---------------------------------------------------------------------------
%% TITLE PAGE WITH AUTHOR DETAILS (pages 1)
%% Comment out this entire block for the anonymous submission PDF.
%% ---------------------------------------------------------------------------
\begin{center}
{\LARGE{Title Page for \textit{The American Mathematical Monthly}}}
\end{center}

\vspace{100px}
\begin{center}

\textbf{\large The continuous-coefficient Jensen equation:} \\
\textbf{\large A note on three vestigial regularity hypotheses} \\

\vspace{30px}

%% Author 1 --- fill in before submission of the ``with author details'' PDF.
\mbox{[Author Name]} \\
\mbox{[Institution / Affiliation]} \\
\mbox{[City, Country]} \\
\texttt{\mbox{[author@institution.tld]}}

\end{center}

\vspace{20px}
\begin{center}
\textbf{Author biography} (2--5 sentences; required for Articles):

\mbox{[Brief biographical note: current position, institution, research interests,}
\mbox{and any relevant background that helps the reader understand the author's}
\mbox{perspective on the problem. Do not include this in the anonymous PDF.]}
\end{center}

\pagebreak
%% ---------------------------------------------------------------------------
%% ANONYMOUS MANUSCRIPT BODY (double-blind: no author name below this line)
%% ---------------------------------------------------------------------------

\title{The continuous-coefficient Jensen equation:\\
       A note on three vestigial regularity hypotheses}

%% Running head: abbreviated title, ≤50 characters.
\markright{Continuous-coefficient Jensen equation}

%% DO NOT fill in author names until the PROVISIONAL ACCEPT letter.
%% Monthly submissions are DOUBLE BLIND.
\author{Double Blind No Name}

\maketitle

%% ---------------------------------------------------------------------------
%% ABSTRACT
%% Guidelines: ≤250 words; entice the reader; minimize notation;
%%             highlight concepts, not mechanics.
%% ---------------------------------------------------------------------------
\begin{abstract}
Jensen's inequality says that for a concave function~$G$ and a random
variable~$\xi$ on an interval, $\E[G(\xi)]\le G(\E[\xi])$. A natural
question: if equality holds for every two-point distribution---if
$p\,G(u_1)+(1-p)\,G(u_2)=G(p\,u_1+(1-p)\,u_2)$ for all $u_1,u_2$ in
the interval and all weights $p\in[0,1]$---must $G$ be affine?

The answer is yes. The proof is one line: substitute $u_1=M$,
$u_2=0$, $p=v/M$ to read $G(v)$ directly off the chord through the
endpoints. No regularity hypothesis on~$G$---continuity, measurability,
monotonicity, or boundedness---is needed.

This contrasts with a celebrated 80-year story. The $p=\tfrac{1}{2}$
specialization is equivalent to Cauchy's additive equation and admits
wildly pathological non-affine solutions without regularity. The Polish
school of the 1920s cataloged which tameness conditions---continuity,
measurability, boundedness on a set of positive measure---suffice to
rule out the pathology. These conditions became standard cargo in
functional-equation arguments, carried reflexively even where the full
continuous-weight equation is in play and no such cargo is needed.

We make explicit which classical hypotheses become vestigial once all
weights $p\in[0,1]$ are tested (rather than rational weights only),
explain structurally why (the pathology lives at irrational~$p$,
precisely where the continuous equation tests~$G$), and identify the
recurring trap: derivations in expected-utility theory, entropy
axiomatics, and surrogate calibration that saturate Jensen across the
full two-point family each produce this equation and each benefit from
the one-line closure.
\end{abstract}

\vspace{8px}

\noindent\textbf{Keywords:} Jensen equation, Cauchy equation, Hamel basis,
functional equation, regularity hypothesis

\vspace{4px}

\noindent\textbf{MSC 2020:} 39B22, 39B05

\vspace{8px}

%% =============================================================================
\section{Introduction.}
\label{sec:intro}

One of the most serviceable inequalities in mathematics---Jensen's
inequality, $\E[G(\xi)]\le G(\E[\xi])$ for concave~$G$---raises an
immediate follow-up: when does equality hold for \emph{every} two-point
distribution supported on a given interval? The equation expressing
universal two-point equality,
\begin{equation}
p\,G(u_1)+(1-p)\,G(u_2) \;=\; G\bigl(p\,u_1+(1-p)\,u_2\bigr)
\qquad (u_1,u_2\in I,\ p\in[0,1]),
\tag{$\star$}\label{eq:star}
\end{equation}
imposed on a function $G\colon I\to\R$ on a real interval $I$, has
at least three faces in classical analysis. Those faces have been steadily
conflated for over a century, with consequences for both the foundational
theory and the modern applied literature where the equation surfaces.

\subsection{Three faces of the equation.}

The \emph{discrete-coefficient} Jensen equation,
\begin{equation}
G\!\left(\tfrac{u_1+u_2}{2}\right) \;=\; \tfrac{G(u_1)+G(u_2)}{2},
\tag{$J_2$}\label{eq:J2}
\end{equation}
is the $p=\tfrac{1}{2}$ specialization of~\eqref{eq:star}. Setting
$f(x):=G(x)-G(0)$, the equation~\eqref{eq:J2} is equivalent on a
translate of~$I$ to Cauchy's additive equation $f(x+y)=f(x)+f(y)$.
Cauchy's equation inherits a full apparatus of pathological solutions:
without measurability, monotonicity, boundedness on a set of positive
Lebesgue measure, or another regularity hypothesis, the equation admits
non-affine solutions constructed via a \emph{Hamel basis}---a basis
of~$\R$ as a $\Q$-vector space, whose existence requires the axiom of
choice. The classical pathological-solution apparatus is the work of
Hamel~\cite{Hamel1905}; the regularity hypotheses sufficient to recover
affineness were cataloged by Cauchy~\cite{Cauchy1821},
Darboux~\cite{Darboux1875}, Sierpi\'nski~\cite{Sierpinski1920},
Steinhaus~\cite{Steinhaus1920}, and Ostrowski~\cite{Ostrowski1929}.

The \emph{rational-coefficient} Jensen equation,
\begin{multline}
p\,G(u_1)+(1-p)\,G(u_2)
  \;=\; G\bigl(p\,u_1+(1-p)\,u_2\bigr) \\
  (u_1,u_2\in I,\quad p\in[0,1]\cap\Q),
\tag{$J_\Q$}\label{eq:JQ}
\end{multline}
strengthens~\eqref{eq:J2} but does not escape its pathology. Cauchy
additivity entails $\Q$-homogeneity---$f(qx)=qf(x)$ for every $q\in\Q$
and $x\in\R$---by an elementary induction on positive integers, sign
reversal for negatives, and the inverse-of-multiplication argument for
$1/n$. Combining additivity with $\Q$-homogeneity
returns~\eqref{eq:JQ} for~$f$, hence for~$G$. So \eqref{eq:JQ},
\eqref{eq:J2}, and Cauchy's equation share the same solution class up to
constants, and all three inherit the Hamel-basis pathology.

The \emph{continuous-coefficient} form~\eqref{eq:star}, by contrast, is
fundamentally different. The strengthening from $\Q$ to $\R$ in the
range of~$p$ is \emph{genuine}. As we will see, \eqref{eq:star} does not
admit any Hamel pathology: the pathology cannot survive at irrational~$p$,
and the continuous-coefficient equation tests~$G$ at every irrational.
This is the subject of the paper.

\subsection{The main result.}

\begin{theorem}\label{thm:main}
Let $M>0$ and let $G\colon[0,M]\to\R$ satisfy~\eqref{eq:star} for every
$u_1,u_2\in[0,M]$ and every $p\in[0,1]$. Then $G$ is affine on $[0,M]$:
\[
G(v) \;=\; \frac{G(M)-G(0)}{M}\,v + G(0), \quad v\in[0,M].
\]
No regularity hypothesis on $G$ (continuity, measurability,
monotonicity, boundedness) is required.
\end{theorem}

The proof is one line: a substitution into~\eqref{eq:star} that pins
$G(v)$ along the chord through the endpoints $(0,G(0))$ and $(M,G(M))$.
We give it in Section~\ref{sec:proof} and use the remainder of the paper
to articulate what it means.

\subsection{The historical detour and the structural moral.}
\label{ssec:detour}

Theorem~\ref{thm:main} is folklore (cf.~\cite[\S2.1]{Aczel1966},
\cite[Ch.~13]{Kuczma2009}). It is \emph{not} a new mathematical result.
The contribution of the present paper is in three parts:
\begin{itemize}
\item \emph{The explicit dictionary} (Section~\ref{sec:dictionary}) of
  three regularity hypotheses on~$G$---continuity, measurability,
  boundedness---that are required to force affineness
  for~\eqref{eq:J2} but become \emph{vestigial}
  under~\eqref{eq:star}, together with the structural mechanism (the
  Hamel pathology lives at irrational~$p$, exactly where
  \eqref{eq:star}'s continuous coefficient closes the door).

\item \emph{The strict-minimum hypothesis}
  (Section~\ref{ssec:weak})---the
  equation~\eqref{eq:star} at a single chord configuration already
  suffices.

\item \emph{The recurrence pattern}
  (Section~\ref{sec:recurrence})---the
  trap of invoking unnecessary Cauchy--Hamel regularity recurs whenever
  a derivation pushes Jensen's inequality to saturation across a wide
  class of two-point distributions. We survey three example application
  areas where this pattern arises.
\end{itemize}

The 80-year detour through Hamel and the Polish school answered a
perfectly natural question---\emph{what regularity hypothesis suffices
for \eqref{eq:J2}$\Rightarrow$affine?}---but a \emph{different}
question from the one that matters for many modern applications. The
applied derivation does not produce \eqref{eq:J2} or \eqref{eq:JQ};
it produces \eqref{eq:star} as a saturated identity, and the one-line
substitution of Section~\ref{sec:proof} is the right closure tool. The
dictionary of Section~\ref{sec:dictionary} makes explicit which classical
regularity hypotheses become vestigial in this transition.

\paragraph{Terminology.}
Throughout the paper we refer to the substitution $u_1=M$, $u_2=0$,
$p=v/M$ in~\eqref{eq:star}---by which $G$ is pinned along the chord
through the endpoints $(0,G(0))$ and $(M,G(M))$---as the
\emph{chord substitution}. The substitution is standard folklore in
the functional-equations literature
(cf.~\cite[\S2.1]{Aczel1966}, \cite[Ch.~13]{Kuczma2009}); the
descriptive label is our convenience, adopted for brevity in repeated
reference rather than as a claim of established terminology.

%% =============================================================================
\section{The result and its proof.}
\label{sec:proof}

\begin{proof}[Proof of Theorem~\ref{thm:main}]
Fix $v\in[0,M]$. Set
\[
u_1 := M, \qquad u_2 := 0, \qquad p := \frac{v}{M}.
\]
The constraints of~\eqref{eq:star} are satisfied: $u_1=M\in[0,M]$,
$u_2=0\in[0,M]$, and $p=v/M\in[0,1]$ since $v\in[0,M]$.
Substituting into~\eqref{eq:star},
\[
\frac{v}{M}\,G(M)+\left(1-\frac{v}{M}\right)G(0)
\;=\;
G\!\left(\frac{v}{M}\cdot M+\left(1-\frac{v}{M}\right)\cdot 0\right)
\;=\; G(v).
\]
Rearranging yields $G(v)=G(0)+\bigl(G(M)-G(0)\bigr)\,v/M$.
\end{proof}

The chord substitution pins $G(v)$ along the straight chord from
$(0,G(0))$ to $(M,G(M))$ for every $v\in[0,M]$. Since this determines
$G(v)$ at every~$v$, the function~$G$ is fully determined by its values
at the two endpoints---it is exactly the affine function joining them.

\begin{corollary}\label{cor:regularity}
Under the hypotheses of Theorem~\ref{thm:main}, $G$ is in particular
continuous, monotone, locally Lipschitz, absolutely continuous, and
Lebesgue measurable on $[0,M]$.
\end{corollary}

\begin{proof}
Every affine function $G(v)=av+b$ on a bounded interval is differentiable
with constant derivative $G'=a$, hence continuous; non-decreasing if
$a\ge0$ and non-increasing if $a\le0$; Lipschitz with constant~$|a|$;
absolutely continuous; and Borel-measurable, hence Lebesgue-measurable.
\end{proof}

The order of inference matters and is worth stating explicitly. Under
\eqref{eq:star}, the regularity properties of~$G$ are the
\emph{conclusion}, not the hypothesis. Each regularity hypothesis one
might be tempted to assume---continuity, monotonicity, boundedness,
measurability---is automatically satisfied by any solution
of~\eqref{eq:star}. The continuous-coefficient equation supplies its own
regularity from within; no external regularity hypothesis is needed.

This is what makes the continuous-coefficient case structurally different
from the discrete-coefficient case~\eqref{eq:J2}, where some regularity
hypothesis \emph{is} required to rule out the Hamel pathology. We now
examine the difference in detail.

%% =============================================================================
\section{The dictionary of vestigial regularity hypotheses.}
\label{sec:dictionary}

\subsection{Six hypotheses and their fates.}
\label{ssec:dict-table}

We place side-by-side the discrete-coefficient equation~\eqref{eq:J2}
and the continuous-coefficient equation~\eqref{eq:star}, and ask which
regularity hypotheses on~$G$ are \emph{required} to force affineness
for each. The first column of Table~\ref{tab:dictionary} names the
hypothesis; the second records the requirement for~\eqref{eq:J2} with
the standard historical attribution; the third records what becomes of
the requirement under~\eqref{eq:star}.

\begin{table}[ht]
\renewcommand{\arraystretch}{1.4}
\centering
\begin{tabularx}{\linewidth}{@{} >{\raggedright\arraybackslash}X |
  >{\raggedright\arraybackslash}X |
  >{\raggedright\arraybackslash}X @{}}
\toprule
\textbf{Hypothesis on $G$} &
\textbf{For \eqref{eq:J2}$\Rightarrow$affine?} &
\textbf{For \eqref{eq:star}$\Rightarrow$affine?} \\
\midrule
\textbf{Continuity on $I$}: $G$ is continuous at every $v\in I$.
  & Yes, suffices (Cauchy~\cite{Cauchy1821}).
  & \textbf{No}---it is the \emph{conclusion} (Corollary~\ref{cor:regularity}). \\
\midrule
\textbf{Measurability on $I$}: $G$ is Lebesgue- or Borel-measurable on~$I$.
  & Yes, suffices (Sierpi\'nski~\cite{Sierpinski1920}).
  & \textbf{No}---same. \\
\midrule
\textbf{Monotonicity on $I$}: $G$ is non-decreasing or non-increasing on~$I$.
  & Yes, suffices (Darboux~\cite{Darboux1875}).
  & \textbf{No}---same. \\
\midrule
\textbf{Boundedness on a set of positive measure}: $G$ is bounded on some
  Lebesgue-measurable $E\subseteq I$ with $|E|>0$.
  & Yes, suffices (Steinhaus~\cite{Steinhaus1920}).
  & \textbf{No}---same. \\
\midrule
\textbf{Boundedness on $I$}: $G$ is bounded on all of~$I$.
  & Yes, suffices (Ostrowski~\cite{Ostrowski1929}).
  & \textbf{No}---same. \\
\midrule
\textbf{None}: no regularity hypothesis at all.
  & \textbf{Insufficient}---Hamel-basis pathology~\cite{Hamel1905}.
  & \textbf{Sufficient}---Theorem~\ref{thm:main}. \\
\bottomrule
\end{tabularx}
\caption{Regularity hypotheses on $G$. Column~2 records which hypotheses
are classically required to force affineness for the
discrete-coefficient equation~\eqref{eq:J2}. Column~3 records their
status under the continuous-coefficient equation~\eqref{eq:star}.}
\label{tab:dictionary}
\end{table}

The bottom row is the point of the table. \textbf{No regularity
hypothesis} suffices for~\eqref{eq:J2}: there exists a non-affine,
non-measurable, unbounded, additive function on~$\R$---the classical
Hamel-basis pathology~\cite{Hamel1905}---that supplies a non-affine
solution of~\eqref{eq:J2}. Conversely, \textbf{no regularity hypothesis}
is \emph{needed} for~\eqref{eq:star}: Theorem~\ref{thm:main}'s chord
substitution closes the proof at the level of the bare functional
equation.

The five non-bottom rows tell a complementary story. For~\eqref{eq:J2},
\emph{some} hypothesis from these five is required; otherwise the Hamel
pathology survives. For~\eqref{eq:star}, \emph{none} of these is
required---and yet each is automatically satisfied by any solution
(Corollary~\ref{cor:regularity}). Each hypothesis is, in the
dictionary's terminology, \emph{vestigial}: classically required
for~\eqref{eq:J2}, classically expected for~\eqref{eq:star}, but in
fact unnecessary for~\eqref{eq:star}.

\subsection{The structural mechanism.}
\label{ssec:mechanism}

Why does the Hamel pathology fail to apply to~\eqref{eq:star}? The
mechanism is direct and illuminating.

A Hamel-basis pathological solution $G$ of Cauchy's equation on~$\R$ is
$\Q$-linear: $G(0)=0$ and $G(qx)=q\,G(x)$ for every $q\in\Q$ and
$x\in\R$. By design, $G$ is \emph{not} $\R$-linear: there exists an
irrational $r\in\R$ and an $x\in\R$ such that $G(rx)\ne r\,G(x)$.

The chord identity in~\eqref{eq:star} at the configuration $u_1=M$,
$u_2=0$ reads
\begin{equation}
p\,G(M)+(1-p)\,G(0) \;=\; G(p\,M).
\tag{$\star_0$}\label{eq:star0}
\end{equation}
For a Hamel-pathological $G$ (so $G(0)=0$), \eqref{eq:star0} becomes
$p\,G(M)=G(p\,M)$---exactly the $\R$-linearity assertion for~$G$ at the
pair $(M,p)$. By $\Q$-linearity of~$G$, this holds for every $p\in\Q$.
By the failure of $\R$-linearity, it \emph{fails} at some irrational~$p$.
So the Hamel-pathological $G$ satisfies~\eqref{eq:JQ} (every rational
coefficient is OK) but violates~\eqref{eq:star} at every irrational~$p$
where the $\R$-linearity defect appears.

The structural point is one sentence:

\begin{quote}
\textbf{The Hamel pathology lives at irrational $p$; the
continuous-coefficient equation~\eqref{eq:star} forecloses on it there.}
\end{quote}

The 80-year detour through Hamel and the Polish school---Cauchy 1821 to
Ostrowski 1929---was an investigation of what regularity hypothesis can
substitute for the \emph{unavailable} constraint at irrational~$p$,
when only the rational-coefficient version of the equation is in play.
With the continuous-coefficient version, the irrational-coefficient
constraint is \emph{available}, and the entire dictionary collapses.

%% =============================================================================
\section{Variants and limits.}
\label{sec:variants}

\subsection{The strict-minimum hypothesis.}
\label{ssec:weak}

The proof of Theorem~\ref{thm:main} uses~\eqref{eq:star} at exactly one
configuration: $u_1=M$, $u_2=0$, with $p\in[0,1]$ ranging freely. The
full force of~\eqref{eq:star}---for every pair $(u_1,u_2)$---is not
consumed. We record the strict minimum.

\begin{theorem}\label{thm:weak}
Let $M>0$ and let $G\colon[0,M]\to\R$ satisfy
\[
p\,G(M)+(1-p)\,G(0) \;=\; G(p\,M) \qquad \text{for all }p\in[0,1].
\]
Then $G(v)=G(0)+(G(M)-G(0))\,v/M$ on $[0,M]$.
\end{theorem}

\begin{proof}
Set $p:=v/M$; then $p\in[0,1]$ since $v\in[0,M]$. The hypothesis gives
\[
p\,G(M)+(1-p)\,G(0) \;=\; G(p\,M) \;=\; G(v),
\]
so $G(v)=G(0)+p\bigl(G(M)-G(0)\bigr)=G(0)+\bigl(G(M)-G(0)\bigr)v/M$.
\end{proof}

Theorem~\ref{thm:weak} is \emph{in principle} weaker than
Theorem~\ref{thm:main}: the one-chord hypothesis does not a priori imply
the full~\eqref{eq:star}. Under the conclusion (affineness) both hold;
in practice an author who derives~\eqref{eq:star} from a richer setup
(e.g., a Jensen-equality identity in a Bayes-risk computation;
see Section~\ref{sec:recurrence}) typically obtains~\eqref{eq:star} in
full force. The narrower Theorem~\ref{thm:weak} is the right reference
when one wants to know exactly how much of the equation the proof
actually consumes.

\subsection{Convex domains in higher dimensions.}
\label{ssec:higher}

The chord substitution extends to any convex subset of a real vector
space.

\begin{theorem}\label{thm:higher}
Let $V$ be a real vector space, let $C\subseteq V$ be a convex set, and
let $G\colon C\to\R$ satisfy
\[
p\,G(x_1)+(1-p)\,G(x_2) \;=\;
G\bigl(p\,x_1+(1-p)\,x_2\bigr) \qquad (x_1,x_2\in C,\ p\in[0,1]).
\]
Then there is a linear functional $a\colon\mathrm{span}(C-C)\to\R$ and a
constant $b\in\R$ such that $G(x)=a(x-x_0)+b$ on~$C$ for any fixed
$x_0\in C$. Equivalently, $G$ is affine on~$C$.
\end{theorem}

\begin{proof}[Proof sketch]
Apply Theorem~\ref{thm:main} along every chord in~$C$: for each
direction $v$ with $x_0+v\in C$, the map
$\lambda\mapsto G(x_0+\lambda v)$ is affine in $\lambda\in[0,1]$. Hence
$G(x_0+v)-G(x_0)$ depends linearly on~$v$ in each direction; an
inductive argument on the dimension of the affine span of the points
involved promotes this to a global linear functional~$a$ on
$\mathrm{span}(C-C)$. The bookkeeping is standard and is given
in~\cite[Ch.~13]{AczelDhombres1989}.
\end{proof}

The structurally essential step is the one-dimensional chord substitution
of Theorem~\ref{thm:main}. The higher-dimensional lifting is bookkeeping.

\subsection{The rational-coefficient version retains the pathology.}
\label{ssec:JQ-pathology}

The strengthening from rational to continuous coefficients
in~\eqref{eq:star} is doing genuine work: the rational-coefficient
version~\eqref{eq:JQ} retains the Hamel pathology.

\begin{proposition}[Folklore; cf.~{\cite[Ch.~5]{Kuczma2009}}]%
\label{prop:JQ-pathology}
There exists $G\colon[0,1]\to\R$ satisfying~\eqref{eq:JQ}
(hence~\eqref{eq:J2}) on $[0,1]$ that is not affine.
\end{proposition}

\begin{proof}[Construction]
Choose a Hamel basis $H$ of~$\R$ over~$\Q$ containing~$1$; such a basis
exists by Zorn's lemma applied to the partially ordered set of
$\Q$-linearly independent subsets of~$\R$. Since $H$ has cardinality
$\mathfrak{c}$ and $\Q\cup\{1\}$ is countable, $H$ contains an irrational
$h\in(0,1)$. Define a $\Q$-linear map $\ell\colon\R\to\R$ by setting
$\ell(1)=0$, $\ell(h)=1$, and extending $\Q$-linearly on the remaining
basis elements.

By construction, $\ell$ is $\Q$-linear: $\ell(x+y)=\ell(x)+\ell(y)$
(Cauchy's equation) and $\ell(qx)=q\,\ell(x)$ for all $q\in\Q$
($\Q$-homogeneity). Set $G:=\ell|_{[0,1]}$; then $G$
satisfies~\eqref{eq:JQ} on $[0,1]$ because $\Q$-rational convex
combinations of points of $[0,1]$ remain in $[0,1]$ and $\ell$
respects $\Q$-linear combinations.

Explicit witness that $G$ is not affine: $G(h)=\ell(h)=1\ne0$, while
$G(q)=\ell(q)=q\,\ell(1)=0$ for every rational $q\in[0,1]$. Any
affine map $A\colon[0,1]\to\R$ agreeing with $G$ on $\Q\cap[0,1]$
would satisfy $A(0)=0$ and $A(q)=0$ for all rational $q$---forcing
$A\equiv0$, contradicting $G(h)=1$.
\end{proof}

Proposition~\ref{prop:JQ-pathology} sharpens the dictionary:
\eqref{eq:JQ} on $[0,1]$ admits non-affine solutions; only the
strengthening to the continuous-coefficient~\eqref{eq:star} retires the
regularity-hypothesis requirement.

%% =============================================================================
\section{Where the trap recurs: Jensen saturation as a structural pattern.}
\label{sec:recurrence}

The chord substitution of Section~\ref{sec:proof} is folkloric---both
\cite[\S2.1]{Aczel1966} and \cite[Ch.~13]{Kuczma2009} articulate the
underlying observation, and it is standard in the functional-equations
community. Why, then, does this paper exist?

Because the explicit packaging matters. The functional-equations
community has internalized the chord substitution; the applied
communities that produce the continuous-coefficient equation as a
byproduct of Jensen saturation have not. These vestigial regularity
hypotheses---continuity, measurability,
and boundedness---keep reappearing in applied derivations
of~\eqref{eq:star} outside the functional-equations community. The trap
is not a one-off oversight; it has a structural source. This section
identifies the source and surveys three application areas where it
surfaces.

\subsection{The structural source: saturated Jensen.}
\label{ssec:saturated-jensen}

The recurrence is predictable. Whenever a derivation in any field
arrives at an identity of the form
\begin{equation}
\E[g(\xi)] \;=\; g\bigl(\E[\xi]\bigr)
\qquad (\xi \text{ a random variable on }I,\ g\colon I\to\R),
\tag{$\dagger$}\label{eq:dagger}
\end{equation}
for a class of random variables~$\xi$ wide enough that the marginal
$\E[\xi]$ can be any point of~$I$ and the support of~$\xi$ can be any
two-point subset of~$I$ with any pair of masses $(p,1-p)$,
$p\in[0,1]$, the identity is exactly~\eqref{eq:star} with $g=G$ and $\xi$
supported on $\{u_1,u_2\}$ with masses $\{p,1-p\}$. Such an identity
arises whenever \textbf{Jensen's inequality is pushed to
saturation}---equality is required in $\E[g(\xi)]\le g(\E[\xi])$ for
concave~$g$ across the full two-point family.

Jensen's inequality is the central tool in dozens of fields. In
\emph{most} of these fields, Jensen is used as an inequality with
explicit slack (tracked in convex analysis via biconjugates and Fenchel
duality; in calibration theory via the $\psi$-transform of Bartlett,
Jordan, and McAuliffe~\cite{BJM2006}; in entropy via convexity
arguments). Such derivations never produce~\eqref{eq:star} as a
saturated identity.

In \emph{some} fields, a derivation produces~\eqref{eq:dagger} as a
saturated identity: the slack is forced to zero by a structural property
of the setup (e.g., the class of~$\xi$ is rich enough to exhaust the
inequality's freedom). At that moment~\eqref{eq:star} appears, and the
chord substitution is the natural closure. An author unfamiliar with the
chord substitution will reach for the classical Cauchy--Hamel toolkit and
find that the regularity hypothesis is unnecessary only after the fact.

\subsection{Application: expected-utility representation theorems.}
\label{ssec:utility}

In \textbf{expected-utility theory} in the von Neumann--Morgenstern
tradition~\cite{vNM1944}, the 1944 axiomatization of preferences over
lotteries---sharpened by Herstein and
Milnor~\cite{HersteinMilnor1953}---derives a utility functional
$U$ on a space of lotteries $\mathcal{L}$ that is \emph{linear in
probability}:
\[
U\bigl(p\,L_1+(1-p)\,L_2\bigr)
\;=\; p\,U(L_1)+(1-p)\,U(L_2) \qquad (L_1,L_2\in\mathcal{L},\ p\in[0,1]).
\]
This is exactly~\eqref{eq:star} with $U$ as~$G$ and lotteries as
points. The classical closure of the linearity step uses the
\emph{Archimedean axiom}: for any three lotteries $L_1\succ L_2\succ L_3$,
there exist $p,q\in(0,1)$ such that
$p\,L_1+(1-p)\,L_3\succ L_2\succ q\,L_1+(1-q)\,L_3$.
The Archimedean axiom forces continuity-in-probability of preferences,
which in turn forces continuity of~$U$.

But the chord substitution of Theorem~\ref{thm:main} \emph{also}
closes the linearity step, without any continuity hypothesis. The chord
substitution is the algebraic alternative to the Archimedean axiom.
Even where the Archimedean axiom is invoked (and rightly so, as it
carries other preference-theoretic content), the closure of the
linearity equation~\eqref{eq:star} itself does not require it.

\subsection{Application: Shannon entropy's axiomatic characterization.}
\label{ssec:entropy}

In the \textbf{axiomatic characterization of Shannon entropy} in the
tradition of Khinchin~\cite{Khinchin1957} and
Faddeev~\cite{Faddeev1957}, Shannon's entropy function
$H(p_1,\ldots,p_n)=-\sum_i p_i\log p_i$ is characterized by three
axioms: continuity, maximum at the uniform distribution, and a
\emph{recursivity axiom} that relates the entropy of a joint
distribution to a marginal entropy plus a conditional.

In the modern formalization~\cite[Ch.~5]{Aczel1966}, the recursivity
axiom reduces, in intermediate steps, to a system of functional
equations on auxiliary functions $g,h$ on $[0,1]$. Several of these
intermediate equations are of \eqref{eq:star}-form for the auxiliary
functions, and the closure step uses continuity (the first Khinchin
axiom). The chord substitution of Theorem~\ref{thm:main} is the
alternative closure that does not require continuity as a hypothesis;
in Acz\'el--Dhombres~\cite{AczelDhombres1989} the substitution is part
of the standard toolkit.

\subsection{Application: surrogate calibration on the resolution axis.}
\label{ssec:calibration}

In \textbf{surrogate calibration} in statistical decision theory, the
foundational tradition---Bartlett, Jordan, and
McAuliffe~\cite{BJM2006} and its extensions by Tewari and
Bartlett~\cite{TewariBartlett2007}, Steinwart~\cite{Steinwart2007},
and Reid and Williamson~\cite{ReidWilliamson2010,ReidWilliamson2011}---%
derives its results via convex analysis (biconjugates, $\psi$-transforms,
Fenchel duality) and consequently tracks Jensen as an \emph{inequality}
with explicit slack. These derivations never produce~\eqref{eq:star}
as a saturated identity.

In a recent manuscript of the author~\cite{El2} on the achievable error
floor of partition-based classifiers, the resolution-axis transposition
framing of surrogate calibration produces~\eqref{eq:star} directly: a
two-cell partition computation forces~\eqref{eq:star} on the function~$G$
expressing the partition Bayes risk in terms of an aggregated concave
score, with the cell mass~$p$ ranging freely over $[0,1]$ on an atomless
probability space. In the Lean~4 formalization of~\cite{El2}, the
corresponding lemma was initially declared with a boundedness hypothesis
on~$G$ in deference to the Cauchy--Hamel literature; the proof body then
found that hypothesis unused, by the chord substitution of
Theorem~\ref{thm:main}. This episode is described in the Appendix.

The lesson generalizes. Any derivation that saturates Jensen across a
wide class of two-point distributions should be expected to produce
\eqref{eq:star} and to benefit from Theorem~\ref{thm:main}'s one-line
closure.

\subsection{An invitation.}
\label{ssec:invitation}

The three application areas above are not exhaustive. We invite readers
who have encountered~\eqref{eq:star} in their own work---in mathematical
economics, information theory, statistical decision theory, physics, or
elsewhere---and who have invoked an unnecessary Cauchy--Hamel regularity
hypothesis to extend the catalog. The structural pattern of saturated
Jensen and the chord substitution as its natural closure is, we believe,
recurrent enough to merit explicit articulation.

%% =============================================================================
\section{Concluding remarks.}
\label{sec:conclusion}

Theorem~\ref{thm:main} is, mathematically, one line. The dictionary of
Section~\ref{sec:dictionary} is six rows of a table. The structural
argument of Section~\ref{ssec:mechanism} is one sentence: \emph{the
Hamel pathology lives at irrational~$p$; the continuous-coefficient
equation~\eqref{eq:star} forecloses on it there.} Why does this material
deserve an article in the \textsc{Monthly}?

Two reasons. First, the 80-year detour through Hamel and the Polish
school is a \emph{foundational} result of classical analysis---and the
dictionary of Section~\ref{sec:dictionary} articulates a precise sense in
which the detour answered an adjacent question: not ``what is the minimum
hypothesis for \eqref{eq:star}$\Rightarrow$affine?'' (Theorem~\ref{thm:main}'s
answer: nothing), but ``what is the minimum hypothesis for
\eqref{eq:J2}$\Rightarrow$affine?'' (Polish school's answer: any of a
dictionary of tameness conditions). Both questions are interesting; their
answers are different.

Second, the recurrence pattern of Section~\ref{sec:recurrence} is, to
our knowledge, not articulated explicitly elsewhere in the
functional-equations literature. The pattern---saturated Jensen
produces~\eqref{eq:star}, and the chord substitution closes
it---is recognizable in a number of modern application contexts (utility
theory, entropy, calibration), and the explicit articulation may save
other authors from carrying vestigial regularity hypotheses in their own
work.

The chord substitution itself is the kind of one-line proof that, once
seen, is hard to unsee---but that nonetheless takes some scaffolding to
make visible. The present paper is the scaffolding.

%% =============================================================================
\appendix
\section{How the observation came up: a note on provenance.}
\label{sec:provenance}

\emph{This appendix is non-mathematical and may be skipped by readers
interested only in the mathematics. It describes the route by which the
observation surfaced---through a proof-assistant exercise---and may
interest readers unfamiliar with how mechanized proof can pressure-test
classical heuristics. Readers interested only in the mathematics can
stop at Section~\ref{sec:conclusion}.}

\subsection{The starting position: a defensive hypothesis added on purpose.}
\label{ssec:starting}

The chord substitution surfaced during the formalization of a separate
paper by the present author~\cite{El2} in the \textbf{Lean~4} proof
assistant, paired with the \textbf{Mathlib} mathematical
library. The relevant lemma expressed the step: \emph{if
$G\colon[0,M]\to\R$ satisfies the continuous-coefficient Jensen
equation~\eqref{eq:star}, then $G$ is affine.}

When the lemma was first stated in informal mathematics, the author
added a \textbf{boundedness hypothesis} on~$G$---the image $G([0,M])$
is a bounded subset of~$\R$---for textbook reasons: an internal
adversarial-review round had correctly flagged that, for the
\emph{discrete-coefficient} equation~\eqref{eq:J2}, no affineness
conclusion holds without some regularity (Hamel's pathology;
cf.~Section~\ref{sec:dictionary}). Boundedness on~$I$ is a sufficient
hypothesis by a classical theorem of Steinhaus and Ostrowski
(Table~\ref{tab:dictionary}, row~5). The author expected boundedness
to be load-bearing and declared the Lean lemma accordingly, binding
the boundedness witness to a named variable.

\subsection{The discovery.}
\label{ssec:discovery}

The proof body, when written, did the only natural thing: it set
$u_1:=M$, $u_2:=0$, $p:=v/M$. The equation rearranged in two lines to
the explicit affine formula, and the proof closed---with the boundedness
witness untouched. The proof assistant raised no error (an unused
hypothesis is legal), but the author noticed.

Removing the boundedness witness from the lemma's signature and
re-checking: still valid. The hypothesis the author had carried
defensively was vestigial.

The author's first reaction was suspicion: the literature on Cauchy's
equation is unambiguous that \emph{some} regularity is needed
(Table~\ref{tab:dictionary}). The working assumption had been that the
same held here. The proof assistant said otherwise.

\subsection{What the books say.}
\label{ssec:books}

The canonical references were then consulted:
J.~Acz\'el, \emph{Lectures on Functional Equations and Their
Applications}, \S2.1~\cite{Aczel1966}, and M.~Kuczma, \emph{An
Introduction to the Theory of Functional Equations and Inequalities},
Ch.~13~\cite{Kuczma2009}. Both books treat the continuous-coefficient
Jensen equation~\eqref{eq:star} as a separable problem and
observe---in compressed form, embedded in longer treatments of richer
equations---that \eqref{eq:star} admits the chord-substitution closure
with no regularity hypothesis. The observation is not a discovery of
the present paper: it has been part of the functional-equations toolkit
since at least Acz\'el's 1966 encyclopedia.

What is \emph{not} standard is the \textbf{explicit packaging} of the
observation as a separable lemma with an accompanying dictionary and a
structural explanation of where the pathology lives. Acz\'el and Kuczma
each devote no more than a paragraph to the observation; the dictionary
of Section~\ref{sec:dictionary} and the recurrence pattern of
Section~\ref{sec:recurrence} are, to the author's knowledge, the
contribution of the present paper.

The honest order of events is therefore: \textbf{proof assistant first,
books second.} The author did not find the observation in Acz\'el or
Kuczma and then formalize it; the author formalized the lemma
defensively, the proof assistant showed the hypothesis was unused, and
the literature then confirmed that the observation underlying the closure
is known and properly attributed in the functional-equations community.
This appendix is the record of that order.

\subsection{Two modest morals.}
\label{ssec:morals}

\emph{A defensive hypothesis is not free.} In any application that
needs the lemma, the user must verify the hypothesis. If the hypothesis
is vestigial, that verification is wasted; worse, in a setting where the
hypothesis is not naturally available, the user may abandon a correct
proof route on the mistaken belief that a deeper closure tool is
required. The dictionary of Section~\ref{sec:dictionary} is, in part, an
inventory of regularity hypotheses that---under the continuous-coefficient
form---are exactly such defensive additions.

\emph{Mechanized proof complements mathematical intuition.} The proof
assistant did not discover the chord substitution---Acz\'el and Kuczma
had it long before---but it did force the author to notice that a
hypothesis carried for textbook reasons was unused. That noticing
prompted the literature search that produced the dictionary and the
structural pattern. Without the mechanization, the author would in all
likelihood have written up the lemma with the boundedness hypothesis
attached and never asked the question.

%% =============================================================================
%% Backmatter: Acknowledgments, Disclosure, Funding
%% Note: "Acknowledgments" (American English, no extra 'e').
%% =============================================================================
\section*{Acknowledgments.}

The chord-substitution observation surfaced in the Lean~4 formalization
track of the author's main paper~\cite{El2}. The author thanks the
internal adversarial-review process documented in the source repository
associated with~\cite{El2} for surfacing the over-engineered boundedness
hypothesis whose retirement motivated the present work.

During the preparation of this paper the author used an AI language
model (Anthropic's Claude Opus, accessed via GitHub Copilot) in two
specific capacities: (\emph{i})~adversarial review of successive
drafts, in which the model was prompted in hostile-referee mode to
surface defects in the prose, argument structure, and prior-art
positioning; (\emph{ii})~formatting assistance---\LaTeX{} mechanics and
bibliography construction. All mathematical content, including all
theorems, proofs, and the structural analysis of
Sections~\ref{sec:dictionary} and~\ref{sec:recurrence}, was conceived,
verified, and refined by the author, who takes full responsibility for
the paper.

\section*{Disclosure statement.}

The author reports there are no competing interests to declare.

\section*{Funding.}

This research received no specific grant from any funding agency in the
public, commercial, or not-for-profit sectors.

%% =============================================================================
%% Bibliography (NLM Vancouver style, as required by the Monthly)
%% =============================================================================
\bibliographystyle{vancouver}
\bibliography{refs}

\end{document}
